%!TEX root = ../hutbthesis_main.tex

\cleardoublepage
\sloppy
\begin{thebibliography}{99}
\raggedright
\setlength{\itemsep}{0pt}
\addcontentsline{toc}{chapter}{参考文献}
\bibitem{ref1} 胡子阳,廖胜辉,邬任重,等.融合CBAM注意力机制的人体软组织多物理场快速同步建模[J/OL].山东大学学报(理学版),1-11[2026-04-17].https://link.cnki.net/urlid/37.1389.N.20260113.1404.002.
\bibitem{ref2} 刘代清.手法淋巴引流的皮肤软组织力学建模与淋巴动力学有限元分析[D].电子科技大学,2025.DOI:10.27005/d.cnki.gdzku.2025.002378.
\bibitem{ref3} 张天明.人体运动生物力学建模及仿真研究[D].吉林大学,2024.DOI:10.27162/d.cnki.gjlin.2024.000781.
\bibitem{ref4} 吕晓庆,李旭鸿,童声康.人体生物力学仿真与可视化研究方法进展[C]//中国体育科学学会.第十三届全国体育科学大会论文摘要集——墙报交流(运动生物力学分会).杭州师范大学;,2023:372-374.DOI:10.26914/c.cnkihy.2023.082881.
\bibitem{ref5} 陈卓.基于OpenSim人体肌骨仿真建模的生物力学研究进展[C]//国际班迪联合会(FIB),国际体能协会(ISCA),澳门体能协会(MSCA),中国班迪协会(CBF),中国无线电测向和定向运动协会(CRSOA).2023年首届国际体育科学大会论文集.西北师范大学体育学院;,2023:753-757.DOI:10.26914/c.cnkihy.2023.025181.
\bibitem{ref6} 潘言心.分布式三维力触觉传感器的设计与实现[D].东南大学,2024.DOI:10.27014/d.cnki.gdnau.2024.004197.
\bibitem{ref7} 陈大鹏,丁益,娄隽铖,等.数据驱动的纹理摩擦建模与触觉渲染方法研究[J].仪器仪表学报,2025,46(05):339-351.DOI:10.19650/j.cnki.cjsi.J2513658.
\bibitem{ref8} 陈庚.虚拟纹理多维属性的力触觉建模与再现方法研究[D].南京信息工程大学,2024.DOI:10.27248/d.cnki.gnjqc.2024.001357.
\bibitem{ref9} 王宏润.触觉增强的XR自然手势交互系统研究[D].北京邮电大学,2025.DOI:10.26969/d.cnki.gbydu.2025.003124.
\bibitem{ref10} 刘浩城.面向软组织安全交互的针灸机器人柔性末端控制方法研究[D].广东工业大学,2025.DOI:10.27029/d.cnki.ggdgu.2025.002745.
\bibitem{ref11} 梅欣.基于物体力触觉属性识别的自适应抓取研究[D].扬州大学,2025.DOI:10.27441/d.cnki.gyzdu.2025.000965.
\bibitem{ref12} 马雨前.虚拟手术中基于有限元形变仿真的RPIM无网格建模与碰撞检测研究[D].南昌大学,2024.DOI:10.27232/d.cnki.gnchu.2024.001303.
\bibitem{ref13} 冯上涛.虚拟手术缝合仿真中软组织形变与碰撞检测技术研究[D].南昌大学,2024.DOI:10.27232/d.cnki.gnchu.2024.001092.
\bibitem{ref14} 吕双祺,王伟,贾惟,等.气凝胶复合材料变形行为细观力学建模与仿真[J].科学技术与工程,2025,25(26):11407-11417.
\bibitem{ref15} 李浩达.多体系统非理想运动副接触力学仿真分析[D].哈尔滨工业大学,2025.DOI:10.27061/d.cnki.ghgdu.2025.002875.
\bibitem{ref16} 赵飞.重载履带式爬壁机器人力学建模与仿真[J].机械设计,2025,42(02):148-154.DOI:10.13841/j.cnki.jxsj.2025.02.026.
\bibitem{ref17} Tzimos N ,Parafestas E ,Voutsakelis G , et al.Multimodal Interaction with Haptic Interfaces on 3D Objects in Virtual Reality[J].Electronics,2025,14(20):4035-4035.DOI:10.3390/ELECTRONICS14204035.
\bibitem{ref18} Taghizad H ,Lemaire M ,Massicotte D .Multi-rate real-time simulation: Techniques, models, frameworks, and challenges[J].Array,2026,29100661-100661.DOI:10.1016/J.ARRAY.2025.100661.
\bibitem{ref19} Laukaitis A ,Šareiko A ,Mažeika D .A Dual Digital Twin Framework for Reinforcement Learning: Bridging Webots and MuJoCo with Generative AI and Alignment Strategies[J].Electronics,2025,14(24):4806-4806.DOI:10.3390/ELECTRONICS14244806.
\bibitem{ref20} Islam U N M ,Sultan N ,Saeed A , et al.Enhancing postural stability models: Integrating muscle dynamics and CNS simulation using bond graph modeling[J].Measurement and Control,2026,59(5):682-708.DOI:10.1177/00202940251346099.
\bibitem{ref21} Han Z Z ,Liu X S ,Guo X Y .The orbital dynamics modeling and simulation of asteroid probe based on variational integrators[J].Advances in Space Research,2026,77(5):5986-5995.DOI:10.1016/J.ASR.2025.12.074.
\bibitem{ref22} Mao Z ,Tong S ,Jiang C , et al.How university teachers’ digital literacy influences their innovative ability: a system dynamics theoretical modeling and simulation study[J].Frontiers in Psychology,2026,161665337-1665337.DOI:10.3389/FPSYG.2025.1665337.
\bibitem{ref23} Zhang B ,Xie T Z C .Exploring the complexity of pedestrian dynamics in skiing: A modelling and simulation framework[J].Simulation Modelling Practice and Theory,2026,146103225-103225.DOI:10.1016/J.SIMPAT.2025.103225.
\bibitem{ref24} Kalibala A ,Nada A A ,Ishii H , et al.Dynamic Modeling of a Soft Eversion-Based Growing Robot: Physical Analysis, Simulation, and Experimental Validation.[J].Soft robotics,2026,21695172261442066.DOI:10.1177/21695172261442066.
\bibitem{ref25} Ye M ,Wang J .Dynamic Modeling and Simulation Study of Space Maglev Vibration Isolation Control System[J].Electronics,2026,15(7):1485-1485.DOI:10.3390/ELECTRONICS15071485.
\bibitem{ref26} Zueco C I ,Nicolás L G .Multi-scale elastic contour mapping for deformable object shape control[J].Robotics and Autonomous Systems,2025,194105134-105134.DOI:10.1016/J.ROBOT.2025.105134.
\bibitem{ref27} 戴瑛,吴艾辉,贺鹏飞. 《非线性连续介质力学》课程的知识图谱构建[C]//北京力学会,天津市力学学会,山东省力学学会,河南省力学学会,山西省力学学会. 第二十届北方七省市区力学学会学术会议论文集. 同济大学航空航天与力学学院; ,2025:185-188.DOI:10.26914/c.cnkihy.2025.030601.
\bibitem{ref28} 翟敬梅,章昊. 基于质点-弹簧-阻尼模型的人机接触运动皮肤三维形变仿真[J].清华大学学报(自然科学版),2024,64(10):1706-1716.DOI:10.16511/j.cnki.qhdxxb.2024.26.046.
\bibitem{ref29} 任兆亭,姜学想,李希志,等. 基于滑动平均滤波的高频方波信号注入无传感器控制策略[C]//中国家用电器协会.2025年中国家用电器技术大会论文集(2).青岛海信日立空调系统有限公司; ,2025:45-51.DOI:10.26914/c.cnkihy.2025.072595.
\bibitem{ref30} András B B ,Arik L ,Yannis S , et al.Benchmark of the Physics Engine MuJoCo and Learning-based Parameter Optimization for Contact-rich Assembly Tasks[J].Procedia CIRP,2023,1191059-1064.DOI:10.1016/J.PROCIR.2023.03.149.
\bibitem{ref31} 杨淦华,曾庆军,韩春伟,黄鑫,戴晓强.人机交互遥操作机器人软体手位置跟踪设计与实现[J].工程设计学报,2023,30(2):164-171.
\bibitem{ref32} 方艳红,吴斌,杨正宜.基于球面调和函数的柔性体力触觉建模研究[J].计算机工程与设计,2011,32(9):3091-3094.
\bibitem{ref33} 张宇航,陈雯柏,张佳琪,等.一种面向六自由度机械臂柔顺装配的深度强化学习策略[J].重庆理工大学学报(自然科学),2024,38(12):148-154.
\end{thebibliography}
